\section{Introduction}
\label{sec:intro}

The models used throughout this thesis describe populations that are small enough
for chance to matter. A handful of bacteria deposited in a mammalian host, a few
viral particles inside a single cell: in each case the population is far from the
regime where deterministic rate equations are informative, and the questions that
matter --- will the population establish itself at all? if it survives, how large
will it be? --- are questions about the fluctuations rather than about the mean.
This chapter assembles the mathematical apparatus needed to answer them.

Two threads run through what follows, and it is worth saying at the outset how they
relate.

The first concerns \emph{conditioning on survival}. A subcritical branching process
dies out with probability one, so its unconditional mean decays geometrically and
tells us very little. But if we look only at those realisations that happen still to
be alive at generation $n$, a quite different picture emerges: the conditional
population settles down to a fixed distribution, the \emph{quasi-stationary}
distribution, whose mean is a constant. For the binary Galton--Watson process with
replication probability $p<\tfrac12$ that constant is $1/A(p)$, where
\begin{equation*}
A(p) = \lim_{n\to\infty}\frac{S_n}{(2p)^n},
\end{equation*}
$S_n$ being the probability of survival to generation $n$. Much of this chapter is
concerned with $A(p)$: what it is, how to compute it, and why --- despite a good deal
of effort --- no elementary formula for it has been found.

The second thread concerns \emph{the near-critical regime}. As $p$ approaches
$\tfrac12$ from below the process takes ever longer to die, the quasi-stationary mean
$1/A(p)$ diverges, and every numerical scheme that relies on iterating to
stationarity becomes impractical. Just above $\tfrac12$ the initial count of
replicative individuals comes to weigh heavily on the chance of the process
surviving at all. The behaviour on either side of $p=\tfrac12$ is what makes small
inocula interesting, and it is also what makes them awkward to compute with.

\Cref{sec:gw} sets out the Galton--Watson process and establishes the two facts on
which everything else rests: that extinction is certain at and below criticality,
and that at criticality the expected lifetime is nevertheless infinite.
\Cref{sec:qsd} derives the limiting conditional mean and variance and so introduces
$A(p)$. \Cref{sec:small} looks at what changes when the process starts from a cohort
of $k$ individuals rather than one, and works out the corresponding constant for the
continuous-time birth--death process, where a closed form \emph{is} available.
\Cref{sec:Ap} is devoted to $A(p)$ itself. It gives an infinite-product
representation, an exact series with two-sided bounds that pin the constant between
the continuous-time answer and twice it, the near-critical asymptotic
$A(p)\sim2(1-2p)$, and an account of why the exact function resists closed-form
description --- a question that turns out to be one about the Koenigs linearisation
of the logistic map. \Cref{sec:moc} changes tack and develops the machinery of
generating-function partial differential equations and the method of characteristics
for the absorption models used later in the thesis.

Throughout, results that are standard are presented as such, and the places where
the argument is heuristic or the numerics are doing the real work are marked
explicitly. Where a claim can be checked numerically, it has been.
